\chapter{Trismus}

\section*{Definition}
Trismus is a sustained, involuntary tonic contraction of the masticatory muscles resulting in restricted mouth opening, typically defined as maximal interincisal distance less than 35~mm. It must be distinguished from temporomandibular joint ankylosis, which involves mechanical rather than muscular restriction.

\section*{Pathophysiology}
Trismus arises from sustained hyperactivity of the muscles of mastication -- the masseter, temporalis, and medial pterygoid -- often driven by afferent trigeminal nerve irritation. Spasm is triggered by local inflammation, infection, or nociceptive input that lowers the firing threshold of the alpha motor neurons supplying the mandibular division (V3). Chronic cases may involve scarring of the muscle fibers themselves.

\section*{Causes}
\begin{itemize}
  \item Odontogenic infection: Pericoronitis (especially mandibular third molars), periapical collection of pus, dental caries extending to the submandibular or masticator space
  \item Orofacial infection: Ludwigs angina, parapharyngeal collection of pus, tonsillar collection of pus, parotitis
  \item Trauma: Mandibular fracture, temporomandibular joint injury, post-surgical (third molar extraction, orthognathic surgery, tonsillectomy)
  \item Neoplastic: Oral cavity tumors, nasopharyngeal carcinoma, parotid tumors infiltrating the pterygoid musculature
  \item Medication-induced: Neuroleptic malignant syndrome, acute dystonic reaction to antipsychotics (metoclopramide, haloperidol), phenytoin toxicity
  \item Central nervous system: Tetanus, status epilepticus, cerebral palsy, brainstem stroke, strychnine poisoning
  \item Other causes: Temporomandibular disorder, radiation-induced scarring (post-radiotherapy for head/neck cancer), scleroderma, progressive systemic sclerosis
\end{itemize}

\section*{When to See a Doctor}
See your doctor if:
\begin{itemize}
  \item Inability to open mouth beyond 20~mm or complete inability to speak or eat
  \item Trismus accompanied by fever, drooling, neck swelling, or high-pitched breathing sound (concern for deep neck space infection)
  \item Recent head/neck trauma or surgery with new-onset restricted opening
  \item Rapid onset over hours or trismus associated with muscle rigidity elsewhere (tetanus, neuroleptic malignant syndrome)
  \item Known malignancy or history of head/neck radiation
\end{itemize}

\section*{Related Symptoms}
\begin{itemize}
  \item Jaw pain
  \item Masticatory muscle tenderness
  \item Bruxism
  \item Facial swelling
  \item difficulty swallowing
  \item Referred otalgia
  \item Headache
  \item Locking or clicking of the temporomandibular joint
\end{itemize}
\textbf{Important:} Trismus is a diagnostic clue for deep space infection when accompanied by fever and neck swelling, and for tetanus when generalized rigidity or risus sardonicus is also present.

\section*{Self-Care / Home Management}
\begin{itemize}
  \item Apply moist heat to the masseter region for 15--20 minutes several times daily to promote muscle relaxation.
  \item Perform gentle, passive jaw-opening exercises within pain tolerance using stacked tongue depressors or a commercial Therabite device.
  \item Consume a soft or liquid diet to reduce masticatory demand; avoid hard, chewy, or crunchy foods.
  \item Maintain oral hygiene with a soft-bristled toothbrush to prevent secondary gingival inflammation.
\end{itemize}

\section*{How Your Doctor Will Evaluate You}
\begin{itemize}
  \item Measure maximal interincisal distance with a ruler to quantify severity and track progression.
  \item Inspect the oropharynx, dentition, and gingivae for swelling, purulence, or caries; palpate the submandibular and parotid regions.
  \item Obtain panoramic radiograph (Panorex) for odontogenic sources; CT with contrast of the face and neck for suspected deep space infection or collection of pus.
  \item Evaluate for systemic causes with serum calcium (tetany), creatine kinase, and review of medication list for dystonic agents.
\end{itemize}

\section*{Treatment Options}
\begin{itemize}
  \item Treat underlying odontogenic or oropharyngeal infection with antibiotics (amoxicillin/clavulanate or clindamycin) and surgical drainage of collection of pus where indicated.
  \item Administer muscle relaxants (diazepam, cyclobenzaprine, baclofen) for spasm reduction; botulinum toxin A injection into masseter and temporalis for refractory cases.
  \item Discontinue or adjust offending medications if drug-induced dystonia is suspected; diphenhydramine or benztropine may reverse acute dystonic reactions.
  \item Institute physical therapy including passive stretching, moist heat, and jaw range-of-motion exercises; severe fibrotic cases may require surgical release (coronoidectomy).
\end{itemize}

\clearpage
